\section{Local projection and the perfect-matching bound}
\label{sec:projection-analysis}

We now implement the fine-ledger projection and prove the progress and
cost bounds for the phase algorithm of
Section~\ref{sec:matching-state}.  We use its state notation and the
guards in \eqref{match:initial-guards}, and the priced-step interface of
Lemma~\ref{lem:priced-steps}, whose construction is given later.
All estimates in this section with a subscript \(a\) are for a fixed
parameter \(a\).  Constants in an exponent \(O(1/a)\), in contrast,
are absolute.

The threshold tests use the potential-reduction viewpoint of approximate
minimum-ratio circulations developed by Chen et al.\ and its dynamic
threshold formulation by van den Brand, Liu, and
Sidford~\cite{FlowRandomized,BrandLiuSidford2023}. The logarithmic
objective, local exploration and projection guarantees needed here are
proved directly.

\subsection{A local exact projection}

The exchange network of Lemma~\ref{match:exchange-projection} has one vertex
for every bin and row.  Its source and sink will be denoted by
\(s_*,t_*\), to distinguish them from the size exponent in \(S=2^s\).
For an integral parent ledger and a requested floor \(cw\), write
\begin{equation}
\label{proj:deficits}
 b(d)=(cw_d-y_d)_+,\qquad
 u(d)=(y_d-cw_d)_+,\qquad
 \beta=\sum_d b(d),\qquad P=\sum_d p_d .
\end{equation}
The parent floor is one unit higher:
\begin{equation}
\label{proj:parent}
                         y_d+p_d\ge(c+1)w_d .
\end{equation}
In particular, \(\beta\le P\).  A maximum exchange flow, followed when
necessary by its full residual source-reachable bin set, is the exact
information required by Lemma~\ref{match:exchange-projection}.  The following
lemma obtains that information without scanning every bin.

\begin{lemma}[Local exact projection]
\label{lem:projection}
Suppose an integral parent ledger satisfies
\eqref{proj:parent} in a partition state of a subcubic graph on at most
\(S\) physical vertices.  Its allocations, column residuals, weights,
and positive debts support point access and sparse indexed updates.
Assume \(P=O_a(S)\) for the fixed-parameter asymptotic bound.
Then the exact exchange projection can be implemented in
\begin{equation}
\label{proj:time}
                         (P+B)S^{O(1/a)+o_a(1)}
\end{equation}
total time on every path, at all sufficiently large sizes for fixed
sufficiently large parameters.

The implementation either returns the exact exchange-flow changes
and, if needed, the full residual reachable set \(R_0\); reports a
detected stochastic failure; or suspends on a deterministic
prerequisite check.  It uses \(O(\log S)\) priced-step instances.
Each reached instance with valid prerequisites has its conditional
detected-failure allowance \(2^{-\sigma}S^{-60}\).
All updates to a fresh child ledger and all necessary physical scans
are confined to \(O(P)\) represented incidences and bin weight, apart
from the padding and subpolynomial factors in \eqref{proj:time}.
The stated deterministic prerequisites do not cause suspension in
the sufficiently large fixed-parameter range.
\end{lemma}

The construction tests whether the full exchange maximum is at least
an integer threshold.  An optimistic local network supplies an upper
bound on that maximum, and the same explored region with true
capacities supplies a lower bound.  We first construct and analyze
this test, then show how two of its regions recover the full residual
cut.  The proof of Lemma~\ref{lem:projection} will assemble these steps.

\paragraph{Optimistic and true regions.}
Assume \(\beta>0\) for the threshold construction.  Start the
exploration with every \emph{seed} bin satisfying
\(p_d\ge w_d/2\).  Their total weight is at most \(2P\).
The index of positive integral debts has at most \(P\) entries, so
listing these seeds and computing their deficits takes local work.
Every deficit is a seed: if \(y_d<cw_d\), then
\eqref{proj:parent} gives \(p_d>w_d\).

Expanding a bin exposes all its bin-to-row exchange arcs, its true
source and sink capacities, and all row-to-bin arcs and endpoints of
the rows thereby discovered.  Arc identifiers are retained, so an
already represented arc is not inserted a second time.  An
unexpanded bin reached as an endpoint is a \emph{boundary}.  It has
no represented bin-to-row arcs yet and no positive source capacity,
since it is not a seed.  Give it, optimistically, a sink arc of
capacity \(\beta+1\).  Its true spare is available by a point lookup,
and \eqref{proj:parent} implies
\begin{equation}
\label{proj:boundary-spare}
            u(d)=y_d-cw_d\ge w_d-p_d>w_d/2 .
\end{equation}
In particular \(u(d)\ge1\), by integrality.

For every region so obtained, its unlocked optimistic maximum flow
is at least the full true maximum flow.  Indeed, take an integral
path decomposition of a full flow and discard its cycles.  Each
source path starts at a seed.  Retain it until its first boundary,
if any, and then send it directly to the sink.  The retained
interior arcs are represented with their true capacities, and the
total flow is at most \(\beta\), less than every optimistic boundary
sink capacity.  The \emph{explored true graph} has the same represented
arcs, but uses the true sink capacity at every boundary.  It embeds
into the full exchange graph.

\paragraph{Testing a threshold.}
The two local networks just defined have zero lower bounds.
For optimization we use a third network on the represented arcs:
its capacities are scaled, its logarithms have padded arguments,
and some sink arcs acquire permanent positive lower bounds.
These lower bounds will limit the total explored weight.
Lemma~\ref{prelim:small-violation} will convert its exactly
conserved real points into integral feasibility statements, with
the total violation including the return arc.

For an integer \(1\le b\le\beta\), begin a fresh exploration and
multiply all its ordinary capacities by \(16\).  Introduce a return
arc \(t_*\to s_*\) with variable \(z\), and set
\begin{equation}
\label{proj:threshold-network}
                       T=16b,\qquad W=20000(P+1).
\end{equation}
The objective is
\begin{equation}
\label{proj:objective}
 \Phi(F,z)=
 \sum_{e\ {\rm ordinary}}
  \left[\log(F_e-\underline c_e+\zeta_e^-)
       +\log(\overline c_e+\zeta_e^+-F_e)\right]
                         -W\log(T-z).
\end{equation}
Here \(\underline c_e,\overline c_e\) are the current unpadded
bounds, initially with lower bound zero, and
\(\zeta_e^-,\zeta_e^+>0\) are dyadic pads.
At the current point, write \(g=\nabla\Phi\), and let \(\lambda_e\)
be the sum of the absolute derivatives of the individual weighted
logarithmic terms on coordinate \(e\), as in
Lemma~\ref{lem:priced-steps}.  In particular, this is not in general
\(|g_e|\): the two ordinary-arc derivatives can cancel, whereas
on the return arc \(g_z=\lambda_z=W/(T-z)\).
All flows start at zero.  Added arcs also start at zero, and every
subsequent motion is an exact circulation supplied by
Lemma~\ref{lem:priced-steps}; no client operation changes an existing
flow value.
Figure~\ref{fig:local-projection} shows the threshold network and the
transition from an optimistic boundary to a permanently locked sink.
\begin{figure}[!htbp]
  \centering
  \begin{tikzpicture}[
    font=\small, >=Stealth,
    vertex/.style={circle,draw,fill=white,minimum size=7mm,inner sep=1pt},
    arc/.style={->,semithick},
    capacity label/.style={fill=white,inner sep=2pt},
    panel/.style={draw=black!35,rounded corners=3pt,fill=black!2,
      align=center,inner sep=7pt}
  ]
    \fill[blue!4,rounded corners=4pt] (1.8,-.9) rectangle (6.8,1.1);
    \node[anchor=north,text=black!70] at (4.3,-.92)
      {expanded bin and discovered row};
    \node[vertex] (s) at (0,0) {$s_*$};
    \node[vertex] (di) at (3,0) {$d_i$};
    \node[vertex] (r) at (6,0) {$r$};
    \node[vertex,draw=blue!65!black] (db) at (9,0) {$d$};
    \node[vertex] (t) at (12,0) {$t_*$};
    \node[anchor=north,text=blue!65!black] at (9,-.52) {boundary};
    \draw[arc] (s) -- node[capacity label,above] {$16b(d_i)$} (di);
    \draw[arc] (di) to[bend left=20]
      node[capacity label,above] {$16x_e$} (r);
    \draw[arc] (r) to[bend left=20]
      node[capacity label,below] {$16(L+1)$} (di);
    \draw[arc] (r) -- node[capacity label,above] {$16(L+1)$} (db);
    \draw[arc,blue!65!black] (db) --
      node[capacity label,above] {$16(\beta+1)$} (t);
    \draw[arc] (di.north) to[out=60,in=130]
      node[capacity label,above,pos=.55] {true sink capacity $16u(d_i)$} (t.north);
    \draw[arc] (t.south) -- (12,-1.8) --
      node[capacity label,below] {return flow $z$; objective $-W\log(T-z)$, $T=16b$}
      (0,-1.8) -- (s.south);

    \node[panel,text width=3.5cm] (optimistic) at (2.2,-3.6)
      {Unexpanded boundary\\[2pt]
       sink interval $[0,16(\beta+1)]$\\
       true spare $u(d)>w_d/2$};
    \node[panel,text width=3.8cm] (locked) at (9.6,-3.6)
      {Expanded and locked\\[2pt]
       sink interval $[4u(d),16u(d)]$\\
       include the newly exposed arcs};
    \draw[arc] (optimistic.east) --
      node[above,align=center] {cached flow\\reaches $8u(d)$}
      node[below] {$7u(d)<f<10u(d)$} (locked.west);
  \end{tikzpicture}
  \caption{The local threshold network and the boundary transition in
    Lemma~\ref{lem:projection}. Labels on ordinary arcs are scaled upper
    capacities; omitted unpadded lower bounds are zero. The schematic shows
    one expanded incidence and one boundary incidence. Initially all seeds
    $p_d\ge w_d/2$, and therefore all deficit bins, are expanded. A boundary has no
    outgoing bin-to-row arcs and no positive source capacity until expansion.
    Padded bounds are used by the engine. Expansion keeps the current flow
    and installs a permanent lower lock. If the last step succeeds before a
    new trigger is served, its boundary flow still satisfies
    $f<10u(d)<16u(d)$, so the true sink capacity is respected up to padding.}
  \label{fig:local-projection}
\end{figure}


Start each pad at \(4\delta\).  If a cached argument is at most
\(3\delta\), increase its pad by \(4\delta\) and re-anchor that
coordinate.  Explicit changes of bounds also re-anchor it.
Increasing a pad only increases the objective.  The register and
step-size guarantees, exactly as in Lemma~\ref{match:top-maintenance},
keep all ordinary arguments above \(\delta\) before a new step and
above \(\delta/10\) at a refresh after motion.  A serviced side does
not need another increase without further motion or a change of
its bounds.  The repairs can therefore be made on dirty and newly
inserted coordinates, without a scan of the represented graph at
each step.

The loose iteration and operation caps give, on every guarded
prefix, a total sum of pads
\begin{equation}
\label{proj:pad-budget}
             \alpha:=\sum_{e\ {\rm ordinary}}
                            (\zeta_e^-+\zeta_e^+)<S^{-800}.
\end{equation}
One may bound the number of potential repairs using all \(O(S)\)
ordinary arcs of the full exchange graph; locality is not needed
for this preliminary numerical estimate.  Alternatively, the
external-operation cap gives at most a fixed multiple of \(S^{100}\)
pad settings, each of size \(O(\delta)\), which is smaller than the
right side of \eqref{proj:pad-budget} for large \(S\).
There is no intervening infinite repair chain: a repaired argument
is exactly re-anchored, and the current point does not move during
the service.
All ordinary unpadded-bound violations sum to at most \(\alpha\).
Conservation at the source gives
\begin{equation}
\label{proj:return-range}
                     -\alpha\le z\le16\beta+\alpha .
\end{equation}

At each loop test read \(z\) exactly.  Stop successfully when
\(z\ge T-1/4\).  Otherwise the next step cannot cross \(T\):
the return contribution to its length is
\(W|\Delta z|/(T-z)\), and the step's total length is at most
\(q/100\).  Thus its relative return-argument change is at most
\(q/(100W)\).  In particular even the final successful step leaves
\begin{equation}
\label{proj:return-margin}
                              T-z\ge1/8 .
\end{equation}

\paragraph{Boundary locks and explored size.}
Before taking another pricing step, expand any boundary whose
cached sink flow is at least \(8u(d)\).  Replace its unpadded
optimistic sink interval by
\begin{equation}
\label{proj:lock}
                              [\,4u(d),16u(d)\,] .
\end{equation}
Retain this lower lock for the rest of the threshold test, and add
the newly discovered arcs with flow zero.  Pads can be
reinitialized when the sink bounds change.

The transition is interior with a constant margin.  At the start
of a step, an untriggered boundary has actual flow less than
\(8.2u(d)\): apply the register-error inequality to its
padded-zero lower argument and use \(u(d)\ge1\).
During the step its absolute change is at most
\((q/100)(F_e+\zeta_e^-)\), so its resulting flow is less than
\(10u(d)\).  At a first trigger the same cache guarantee gives
actual flow greater than \(7u(d)\).  Hence the new interval
\eqref{proj:lock} is safely interior.  Triggers are checked on
dirty or newly supplied entries using their latest registers.
Re-anchoring without motion cannot invalidate the estimate.
In particular, if success is detected immediately after a step,
before a newly triggered boundary has been serviced, that boundary
still has
\begin{equation}
\label{proj:unserviced-exit}
                              F_e<10u(d)<16u(d).
\end{equation}
Thus a successful exit does not require an uncharged final scan or
expansion.

Every previously locked sink carries at least \(4u(d)\), apart
from its pad violation; each just-triggering sink carries more
than \(7u(d)\).  All other sink flows are nonnegative apart from
their pad violations, and their sum is \(z\).  Equations
\eqref{proj:return-range} and \eqref{proj:boundary-spare} therefore
bound the total weight of nonseed expanded or triggering bins by
\(8P+1\).  Including seeds, the expanded weight is at most
\begin{equation}
\label{proj:region-size}
                              11(P+1).
\end{equation}
This also bounds physical scanning: an expanded bin has \(2w_d-1\)
vertices of degree at most three.  Per unit of expanded weight
there are at most six bin-to-row incidences, at most eighteen
row-to-bin incidences, and at most eighteen boundary occurrences.
Including source and sink arcs, there are fewer than
\(1000(P+1)\) ordinary arcs.  Each is inserted once, the region only
grows, and each sink is locked at most once.  The same bound thus
controls cumulative topological insertions in this test.

\begin{lemma}[Threshold certificate]
\label{proj:threshold}
An unfailed, nonsuspended threshold test cannot report
\(\mathrm{LOW}\) if the unlocked optimistic maximum is at least
\(T\) in scaled units.  Its success exit implies that the explored
true unscaled maximum is at least \(b\).
\end{lemma}

\begin{proof}
First impose the current integral ordinary bounds, including the
locks, and return interval \([0,16\beta]\).
The current point is exactly conserved.  Its ordinary violation
is at most \(\alpha\), and
\eqref{proj:return-range} adds at most \(\alpha\) more for the
return interval.  The total is at most \(2\alpha<1\), not merely
the ordinary-arc error.  Lemma~\ref{prelim:small-violation} supplies
an exactly feasible integral circulation for these bounds.

Starting from its ordinary flow, one can attain the unlocked
optimistic maximum while preserving all sink locks.  If the
present value is smaller, the ordinary residual network contains
a source-to-sink path; erase cycles to make that path simple.
The only additional lower bounds are on incoming sink arcs.
Their reversals leave the sink and cannot occur on a simple
source-to-sink path.  Every augmentation therefore preserves
the locks, even if the starting flow contained cycles.
This constructs, existentially, a feasible comparator
\((F',z')\) with \(z'\ge T\).

The difference from the current point is a nonzero circulation.
For each ordinary logarithmic term, a negative directional
derivative has magnitude at most one, since the comparator is
feasible even without the pads.  There are fewer than
\(2000(P+1)\) such terms.  The return term contributes
\[
                 \frac{W(z'-z)}{T-z}\ge W
\]
positively.  If \(A_+\) and \(A_-\) are the sums of positive and
absolute negative individual contributions, then
\[
 A_+\ge20000(P+1),\qquad A_-<2000(P+1),\qquad
 \frac{g^{\mathsf T}(F'-F,z'-z)}
      {\|\lambda(F'-F,z'-z)\|_1}
       =\frac{A_+-A_-}{A_++A_-}>\frac12 .
 \]
Only a direction is being compared; evaluating the logarithmic
objective at \(z'\ge T\) is neither intended nor necessary.
Lemma~\ref{lem:priced-steps} rules out \(\mathrm{LOW}\).

At a success exit, replace the optimistic boundary capacities by
the true ones and impose return interval \([T,16\beta]\).
Equation~\eqref{proj:unserviced-exit} shows that no unexpanded
boundary violates its true upper capacity, even before trigger
service.  There are no unexpanded deficits.  The ordinary
violation is at most \(\alpha\); the additional return violation
is at most \(1/4+\alpha\).  The total is less than one.
Lemma~\ref{prelim:small-violation} gives a feasible integral scaled
flow of value at least \(T=16b\) in the explored true graph.
Dividing capacities and flow value by \(16\) proves that its
unscaled maximum is at least \(b\).
\end{proof}

Continue until success, \(\mathrm{LOW}\), detected stochastic
failure, or suspension.  On \(\mathrm{LOW}\), verify by a static
integral maximum-flow computation that the unlocked optimistic
unscaled maximum is below \(b\).  On success, verify that the
explored true unscaled maximum is at least \(b\).
A contradictory exact check causes suspension; Lemma~\ref{proj:threshold}
shows that it cannot occur in an unfailed proceeding
computation.  These calls use
Lemmas~\ref{lem:integral-flow} and~\ref{lem:capacity-scaling}
with size padded to \(O(P+B)\).  Capacity digit scaling in
particular avoids a variable polynomial-capacity exponent when
\(P\) is small relative to \(S\).

Thus each verified test decides the threshold for the full graph.
On \(\mathrm{LOW}\), optimistic domination implies that the full
maximum is below \(b\).  On success, the explored true graph embeds
into the full graph and supplies value at least \(b\).
The locked engine network is discarded after the test; the exact
checks use the unlocked optimistic or true capacities, respectively.

\paragraph{Termination and cost of one test.}
The objective range at loop and service states is
\(O_a((P+1)\log S)\).  There are \(O(P+1)\) ordinary terms with
arguments between \(\delta/10\) and a guarded polynomial upper
bound, and the weighted return term has the same total order by
\eqref{proj:return-margin} and \(W=O(P+1)\).
Pad increases never lower the objective.  An arc insertion or a
lock change costs at most \(O(\log S)\) in objective, and there
are only \(O(P+1)\) such touches.  The
\(\Omega(q^2)\) increase per successful priced step consequently
bounds their number by
\begin{equation}
\label{proj:step-count}
                    O_a((P+1)q^{-2}\log S).
\end{equation}
Use, for example, nominal engine lifetime size \(8192(P+1)\),
including every new endpoint and arc.  Point reads of the return
arc, local pad changes, boundary notifications, and the final
static verification are charged.  No global scan is made at a
pricing step.  Lemma~\ref{lem:priced-steps} and its notification
bound now give \eqref{proj:time} for one threshold.  The guard
\(S^{20}\) on local pricing iterations is eventually larger than
\eqref{proj:step-count} because
\(q^{-2}=O_a(S^{200/a})\) and \(P=O_a(S)\).

\paragraph{Recovering the exact flow and the full peel.}
The threshold tests determine the maximum value, but the matching
procedure also needs the exact allocation changes and, when a deficit
remains, the entire residual source-reachable bin set.  We now recover
this information from the small regions explored by the tests.

\begin{lemma}[Recovering the full residual cut]
\label{proj:local-cut}
Suppose \(\beta>0\), and let \(v\) be the full unscaled exchange maximum.
Suppose a true
explored region of maximum \(v\) is available when \(v>0\), and
an unlocked optimistic region of maximum \(v\) is available when
\(v<\beta\).  The union of the available regions, using true
capacities and original arc identifiers, suffices to compute an
exact full exchange flow and, when \(v<\beta\), the full residual
source-reachable bin set.
\end{lemma}

\begin{proof}
Call the supplied true region the success graph and the supplied
optimistic region the failure graph.  Suppose \(v<\beta\), and
let \(Q\) be a minimum source side of the failure graph.
No boundary belongs to \(Q\), because its
sink arc alone has capacity \(\beta+1>v\).  Consequently every
bin of \(Q\) is expanded and all its outgoing bin-to-row arcs of
positive true capacity are represented.  Every represented row
has all its outgoing row-to-bin arcs, and every positive source
arc is present because all deficits are seeded.  Thus all arcs
leaving \(Q\) in the full true graph occur with the same capacity
in the failure graph.  The same set \(Q\) is a full-graph cut
of capacity \(v\).

The true success graph supplies value \(v\) when \(v>0\).
Solve an exact integral maximum flow in the union with true
capacities, using the failure graph alone when \(v=0<\beta\),
or the success graph alone when \(v=\beta>0\).
The union's value is exactly \(v\), by the cut just proved or by
the total source capacity \(\beta\).
Extend its flow by zero to the full network.

When \(v<\beta\), equality of its value with the capacity of
\(Q\) forces every outgoing cut arc to be saturated and every
incoming cut arc to have zero flow.  Neither an unsaturated
forward arc nor a positive-flow reverse arc can leave \(Q\).
All positive-capacity original arcs needed within \(Q\) are
represented, and every needed reverse arc is the reverse of
an arc carrying positive flow in the union.
Residual reachability computed in this local representation is
therefore the full residual source reachability.  Its bin part
is precisely \(R_0\).
\end{proof}

\begin{proof}[Proof of Lemma~\ref{lem:projection}]
List the positive parent debts to find the seeds and compute
\(\beta\), in \(O(P)\) point operations.  If \(\beta=0\), the
unchanged allocation already satisfies the requested floor, so
giving the fresh child an empty debt map finishes the projection.
The parent ledger remains unchanged.

Otherwise binary search, using the verified threshold tests,
determines the full unscaled maximum \(v\).  Retain a true success
graph for \(v\) when \(v>0\).  If \(v<\beta\), retain an optimistic
failure graph for \(v+1\).  Its unlocked maximum is exactly \(v\):
it is below \(v+1\) by the test, at least \(v\) by optimistic
domination, and integral.  At most a constant number of additional
tests retrieves these regions.  Lemma~\ref{proj:local-cut} then
returns the exact flow and, when needed, its full residual bin set.

The union has size \(O(P+1)\).  Its flow support supplies all
allocation changes.  The set \(R_0\) has weight \(O(P)\), either
by \eqref{proj:region-size} or by applying the parent deletion
inequality to the failed floor-\(c\) projection.  Listing its
assigned rows and their affected allocations requires only
that physical weight and bounded row degree.  No scan of an
unexpanded high-volume bin or of an unchanged parent allocation
is necessary.  This completes the local construction.
There are \(O(\log S)\) thresholds and retrieval tests because
\(\beta\le P=O_a(S)\).  The per-test bound \eqref{proj:time},
the final static flow on the union, and these sparse ledger changes
give the asserted total time.  Each reached priced-step instance
has its stated conditional failure allowance; a detected failure
or suspension is returned to the caller immediately.
The numerical and nominal-size prerequisite checks, including
those used above, are justified uniformly on guarded prefixes
in Lemma~\ref{phase:guard-adequacy}.
\end{proof}

\subsection{Debt, structural work, and progress}

The following estimates apply before the first deterministic
interruption, and do not assume that the graph has a perfect
matching.  A detected failure terminates the test; a returned
structural action or projection satisfies its deterministic
invariants.

\begin{lemma}[Phase accounting]
\label{phase:accounting}
In a phase with initial deficit \(K\), threshold \(L\), and target
\(g_*\), every guarded prefix before completion satisfies
\begin{equation}
\label{phase:peel-budget}
 \sum_{\rm peels}w(R_0)\le2L(g_*-1)+S^{-800},
 \qquad
 \#\{\text{blockers}\}\le2C(g_*-1).
\end{equation}
There are \(O(N+Lg_*)\) expansions, the sum of contraction
arities \(|X|\) is \(O(N+Lg_*)\), and the total tick increase is
\(\tau=O_a(N)\).
\end{lemma}

\begin{proof}
Before the terminal \(g_*\)-th gain there are at most \(g_*-1\)
gains whose ledgers are repaired.  Top contractions preserve
total debt, top expansions consume available slack without
other debt, and the positivity service of
Lemma~\ref{match:top-maintenance} adds less than \(S^{-800}\) debt.
By Lemma~\ref{match:deletion}, each repaired gain adds at most \(2L\)
net debt.  By Lemma~\ref{match:exchange-projection}, every projection
peel consumes at least its weight from the top ledger.
Nonnegative remaining debt gives the first bound in
\eqref{phase:peel-budget}.

The number of new blockers at a peel is at most
\(h(R_0)<Cw(R_0)/L\).  Summing this strict inequality and using
integrality gives the second bound in
\eqref{phase:peel-budget}; the tiny positive remainder is less
than one.  In particular, before the first gain of a phase
with \(g_*=1\), no positive-weight peel can occur.

The bottom certificate of Lemma~\ref{match:bottom-surplus} bounds the
weight removed by a gain by \(2L\).
A removed bin has \(2w_d-1\) physical vertices, and at most one
assigned row per removed bin is materialized with it.  Thus the
number of vertices returned to the paired remainder is bounded
by a constant times the total erased weight.  Every expansion
consumes a pair from that remainder.  Its initial population is
at most \(N\), so expansions number \(O(N+Lg_*)\).
Initially there are at most \(N\) bins and each expansion creates
one new bin.  A contraction decreases their number by
\(|X|-1\), and \(|X|\ge2\); hence the sum of \(|X|\) over
contractions is also \(O(N+Lg_*)\).
Expansions and contractions contribute one tick each, and a
repaired gain contributes \(L\), giving
\(\tau=O(N+Lg_*)\).

Finally \(Lg_*=O_a(N)\).  If \(K<100C\), then \(g_*=1\) and
\(L=O_a(N)\).  Otherwise
\[
 Lg_*\le
 \left(\frac{10000CN}{K}+1\right)\frac{K}{100C}
 =O(N).
 \]
This proves the tick bound.
\end{proof}

\begin{lemma}[No false stall without a stochastic failure]
\label{phase:stall}
If the subcubic graph has a perfect matching, a settled phase
state with fewer than \(g_*\) gains admits neither a valid
\(\mathrm{LOW}\) exit nor a permanent lack of progress.
More precisely, when neither a distinct-bin contact nor an
eligible expansion is available, the top objective has a
circulation direction \(d\) with
\[
                      g^{\mathsf T}d
                           >\tfrac12\|\lambda d\|_1 .
 \]
\end{lemma}

\begin{proof}
Let \(i_*\le g_*-1\) be the gains so far.
For this argument only, realize the matching represented by the
partition.  This is an existence comparison, not an instruction
to materialize the graph at every pricing step.
Taking the symmetric difference with a perfect matching gives \(K-i_*\)
vertex-disjoint augmenting paths~\cite{Berge1957}.  At most the blocker count
meet a blocker.

Orient each remaining path from an unmatched endpoint.
Its first bin is free and is entered at its exposed base.
That bin has no other unmatched vertex and no crossing matching
edge.  At a stall there is no edge to a distinct bin.
If the path exits to a row by a nonmatching edge, its next
matching edge is that row's assigned edge, entering another
bin at its base.  This new bin is assigned; its only possible
crossing matching edge is the one just used on entry.
A simple path must therefore next exit it by a nonmatching
edge.  Continuing this argument, the path must eventually
reach \(U\): it avoids blockers, cannot end inside an assigned
bin, and cannot enter another free bin through an assigned
row edge.

Every such entry uses the unique physical base of its bin.
Simplicity and vertex-disjointness consequently imply that
the prefixes through their first \(U\)-contacts use distinct
bins and rows, both within a prefix and across different
prefixes.  Reassign each prefix row to the preceding bin in
its chain.  This is simultaneously injective and frees the
last bin of every prefix.
If \(Z\) denotes all bins with a \(U\)-contact, one can thus
leave simultaneously free a subset of \(Z\) of size
\begin{equation}
\label{phase:free-contacts}
 r\ge K-i_*-2C(g_*-1)\ge K/2 .
\end{equation}

Scale this conceptual assignment by \(L\).
The desired top column capacities are
\[
                         L+p_d^0-Cw_d ,
\]
which are nonnegative because the current top allocation and
positive slack realize them.  Clip overloaded columns by
removing mass.  An assigned column initially carries \(L\),
so its removed mass is at most \(Cw_d\); the total removed is
at most \(Cw(\mathcal D)\le CN\).
The current top allocation proves that all rows can be filled
under these capacities.  Complete the clipped partial
transportation flow by residual source-to-sink augmentations;
one can scale the dyadic data and use integral path
decomposition for this existence argument.  Simple augmenting
paths do not leave the sink, so column sink flows do not
decrease.  The \(r\) initially empty specified columns receive
in total at most the removed mass.
The resulting allocation \(x'\) has nonnegative slacks \(s'\)
and satisfies
\begin{equation}
\label{phase:alternative-slack}
                           s'(Z)\ge Lr-2CN .
\end{equation}

At a stall, the cache-based eligibility guarantee of
Lemma~\ref{match:top-maintenance} gives \(s_d\le5C\) on \(Z\).
Take the difference from the current allocation and slack as
a direction.  Row sums and column equations are unchanged,
so it is a circulation in the top engine graph.
Every negative individual logarithmic directional contribution
has magnitude at most one, since \(x',s'\ge0\).
Their total is at most
\(|E(\mathcal B)|+|\mathcal D|\le4N\).
The positive sum is at least
\[
 \frac{s'(Z)}{5C}-|Z|
 \ge \frac{Lr-2CN}{5C}-N .
 \]
By \eqref{phase:free-contacts} and
\(L\ge10000CN/K\), this is greater than three times \(4N\).
As \(\lambda\) counts the absolute individual contributions,
the claimed ratio is greater than one half.
Lemma~\ref{lem:priced-steps} excludes \(\mathrm{LOW}\) unless
it explicitly reports failure or suspension.  Its positive
objective increase, with the bounded objective accounting
below, also excludes indefinite pricing without another
structural action.
\end{proof}

\subsection{Implementation and total phase cost}
\label{phase:implementation}

We give the representation and cost details needed to sum the
local work.  The charges are to touched records and actual
engine notifications, rather than to a whole graph per event.

\paragraph{Physical structure.}
Maintain physical vertices, bin and row identifiers, mates,
assignments, incident lists, and edge classifications in
dictionaries and linked lists.  Maintain a set of distinct-bin
contacts and, for each bin, a list of contacts to \(U\) with
its first member accessible.  A register refresh changes a
bin's eligibility flag without scanning its contacts.

At a contraction retain a largest old bin's representative.
Move the smaller bins' member lists and the swallowed rows,
updating their physical incidences and affected ledger
records.  Unchanged exterior incidence allocations retain
their values and need only endpoint retargeting when
necessary.  Sum bin data using \(|X|\) point reads per
retained ledger; do not scan the largest bin's members.
A moved old member's bin size at least doubles during its
uninterrupted stay in bins.  Initial membership, new
expansions, and swallowed row entries begin new charges;
erasures and later returns are covered by
Lemma~\ref{phase:accounting}.  Deletions inspect only removed
bins, their assigned rows, and their affected physical
incidences, applying Lemma~\ref{match:deletion}.
Subcubicity therefore bounds all physical structural touches,
including endpoint retargeting, by
\begin{equation}
\label{phase:structural-touches}
                         O_a((S+Lg_*)\log S).
\end{equation}
Assignment flips in the pair primitive require only abstract
assignment updates.  By Lemma~\ref{match:pair-action}, successful
contractions charge their abstract arity, gains charge \(L\)
up to logarithmic factors, and a final abort costs at most
that latter amount.  All bin-pair work is consequently
\begin{equation}
\label{phase:pair-work}
                    O_a((S+Lg_*+L)\polylog S),
\end{equation}
with the fixed reliability factors included.  Materialization
uses Lemma~\ref{match:materialization} and is charged by the same
physical erasure work or by one final phase-boundary scan.

\paragraph{Top values and objective.}
At a structural edit, read and reset only the affected top
coordinates.  The top column residual is available from
\[
                         y_d^0=s_d-p_d^0+Cw_d ;
\]
it does not require a separate global sum of implicitly
moving allocations.  The aggregation, redirection,
restriction, and tiny transfers determine affected slacks
exactly, preserving column equations without rounding.
Compute all final settings of a batch before submitting
deletions, point changes, and insertions to the engine.
Both old and newly anchored descriptors then have positive
arguments, and no pricing call occurs until the completed
topology again satisfies its equations.
Notifications arising during this service can be queued:
there is no implicit primal motion during the service.
Settle all dirty small-argument and eligibility checks
before choosing the next main action.

Use nominal top-engine lifetime size \(SB\), with an explicit
check on the representation.  Only actual topological
insertions and retargetings consume new incidence slots;
repricing, motion, and positivity repairs on retained
coordinates do not.  Equation~\eqref{phase:structural-touches}
makes this nominal bound sufficient eventually.
For fixed \(a\), Lemma~\ref{phase:accounting} gives
\[
 L=O_a(S),\qquad
 \sum_d p_d^0=O_a(Lg_*+1),\qquad
                         s_d\le L+\sum_d p_d^0 .
\]
Thus at pricing states all top arguments lie in a range
whose logarithmic width is \(O_a(\log S)\), and there are
only \(O(S)\) coordinates.  The top objective range is
\(O_a(S\log S)\).  Each affected coordinate at a structural
edit costs at most \(O_a(\log S)\) in objective, so the
cumulative structural loss is
\(O_a((S+Lg_*)\log^2 S)\).
The donor decrease in a tiny positivity transfer has loss
\(O(\delta)\), and its total is already covered by the
tiny-repair count.  Each successful priced step increases
the objective by \(\Omega(q^2)\).  Hence their number is
\begin{equation}
\label{phase:top-step-count}
                             O_a(Sq^{-2}\log^2 S).
\end{equation}
Dirty-entry repairs and eligibility work charge a constant
number of point operations per notified or structurally
touched entry.  The notification bound of
Lemma~\ref{lem:priced-steps} and
\eqref{phase:top-step-count} put the total top work in
\(S^{1+O(1/a)+o_a(1)}\).

\paragraph{Persistent integral ledgers.}
Store allocations by physical incidence and \(y,p\) by bin
in persistent point maps, for example path-copy balanced
trees.  Maintain an index and sum of positive live debts;
deleted keys are cleared, and new bins are initialized
locally.  Retained levels receive their structural edits
separately.  At a projection from integral parent \(j-1\),
retain the parent's roots from \emph{before} a possible
peel, and modify a copy only along the obtained flow support
and deletions.  Include the exterior loads freed by removed
rows.  The new child's debt map is empty.  Apply preserving
deletions separately to coarser ledgers, and let all deeper
fresh ledgers initially share the new child's roots.
This costs local projection-size work and logarithmic
record overhead, not a whole-vector copy at each reset.
Level \(1\) may build entire maps after its global static
flow computation.

The reset periods of Section~\ref{match:reset} are
\(b_j=B^{a-j}\).
Suppose an event selects level \(j>1\).  Even if the event
advances ticks by \(L\), it crosses no \(b_{j-1}\) boundary:
otherwise a coarser reset would have been selected.
Its integral parent last reset to zero debt in this same
\(b_{j-1}\) period interval; every crossing of such a
boundary resets that parent or a coarser ledger.
Its debt rises by at most \(C\) per expansion and \(2L\)
per gain, and never positively for another action.
Since these actions advance ticks by \(1,L\), respectively,
the debt at the triggering event obeys
\begin{equation}
\label{phase:parent-debt}
                              P\le(C+2)b_{j-1}.
\end{equation}
This reasoning uses both the pre-event and post-event
quotients, so a jump over several periods cannot defeat it.
In particular \(P=O_a(S)\) for \(j>1\).

At most \(\lfloor\tau/b_j\rfloor\) events select level \(j\):
each can be charged to a distinct crossed \(b_j\) boundary.
By \(b_{j-1}/b_j=B\), Lemma~\ref{lem:projection}, and
\eqref{phase:parent-debt}, all \(j>1\) projection costs sum to
\begin{equation}
\label{phase:local-reset-sum}
 \sum_{j=2}^a
   \frac{\tau}{b_j}\bigl((C+2)b_{j-1}+B\bigr)
                    S^{O(1/a)+o_a(1)}
       =S^{1+O(1/a)+o_a(1)} .
\end{equation}
The small-instance padding \(B\) is included here.
For level \(1\), the global exact exchange computation has
cost \(S^{1+o(1)}\).  Since
\(b_0=B^a\ge S\) and \(b_1=b_0/B\), its total is at most
\[
                  \frac{\tau B}{S}S^{1+o(1)}
                         =S^{1+O(1/a)+o_a(1)} .
\]
Together with initialization and the preceding charges,
this proves the desired phase cost on completed guarded
prefixes.

\subsection{Guarded prefixes and the PM guarantee}

\begin{lemma}[Eventual adequacy of the guards]
\label{phase:guard-adequacy}
For every sufficiently large fixed \(a\) and fixed
admissible reliability parameter \(\sigma\), all deterministic
guards of the PM construction are adequate for all
sufficiently large \(S\), on every computation path.
The bounds just proved apply also to a prefix ending in a
failed or interrupted attempted operation.  No perfect-matching
promise is needed for these time and adequacy assertions.
\end{lemma}

\begin{proof}
Consider the prefix through a putative first deterministic
interruption.  The elementary initial checks, such as
\[
 S\ge2^{100},\qquad
 100000C,\ B\le S^{1/20},\qquad q\ge S^{-1},
\]
are eventually satisfied at each fixed sufficiently large
\(a\).  The loose \(S^{20}\) iteration caps and the
\(S^{100}\) external-operation cap suffice first to bound
tiny positivity debt and local pads.  A small argument
increased and exactly re-anchored does not require another
increase without motion or a reset of its bounds.  A top
donor is chosen exactly from a bounded-degree row with sum
\(L\); local repair does not create an unbounded chain of
new small donors.  Thus these numerical error estimates
do not assume the sharper running-time bounds that they
are used to prove.

All completed structural actions and peels preserve the
deterministic ledger and matching invariants.  Their debt
payoffs and the bottom certificates prove
Lemma~\ref{phase:accounting} independently of objective
progress.  Include the present structural action if any.
For a coarse numerical bound, one has
\[
 L\le10000CS+1,\qquad g_*\le S,\qquad
 \text{erased weight}\le4Lg_*+1,
\]
and at most \(N+4Lg_*+2\) expansions, allowing the next
one as well.  Every retained integral ledger has acquired
at most
\begin{equation}
\label{phase:crude-debt}
                      C(N+4Lg_*+2)+2Lg_*
\end{equation}
debt between its resets, even without using
\eqref{phase:parent-debt}.  Column loads are at most
\(L|\mathcal A|\).
Together with \(100000C\le S^{1/20}\), these estimates put
all top data, exchange capacities, local return weights
and scaled capacities, and argument upper bounds well
inside \(S^{20}\).
Thus numerical adequacy does not depend on
parameter-dependent rounding constants in the period bound.

At a local threshold test, the current sink consumption
already bounds every lock-trigger expansion by
\eqref{proj:region-size}, before that expansion is
performed.  Ready arguments and post-motion refreshes have
the required lower margin, and explicit edit batches use
positive old and final settings.
At the top, \eqref{phase:structural-touches}, with one
extra physical-size edit allowance, verifies the nominal
lifetime bound \(SB\) at a possible first overflow.
At a local test, \eqref{phase:parent-debt} and its
\(8192(P+1)\) nominal size give the engine condition
\(R+B\le S^{1+2/a}\) eventually.  The same is true for
the top nominal size.

The structural and objective estimates
\eqref{proj:step-count} and \eqref{phase:top-step-count}
bound all completed iterations.  Add at most one final
\(\mathrm{LOW}\), failed, or interrupted attempt.
Because \(q^{-2}=O_a(S^{200/a})\), these bounds remain
below \(S^{20}\) for fixed sufficiently large \(a\) and
large \(S\).
Up to a hypothetical first external-operation cap,
the engine's notification bound charges dirty coordinates
to priced motion.  Client repairs add only constantly
many point operations per dirty trigger, besides the
structural edits, and have no no-motion repair chain.
Consequently even this cap-reaching prefix has only
\(S^{1+O(1/a)+o_a(1)}\) external work, eventually smaller
than \(S^{100}\).  The priced-step lemma's own eventual
prerequisite adequacy applies within these hypotheses.

These arguments are uniform over successful verified
states and rejected stochastic attempts.  The engine
lemma bounds attempted operations on every path; its
construction uses verified structures and capped failed
attempts, not a good-event-only time estimate.
The pair primitive likewise either returns a valid
action within its charged bound or terminates negatively.
A detected stochastic failure promptly ends the PM test,
whereas an inadequate deterministic prerequisite suspends
the branch.  The latter has just been excluded in the
fixed-parameter asymptotic range.  All local scans and
updates, including a final pair abort or incomplete
attempt, are charged by their represented work.
This proves the claim.
\end{proof}

\begin{theorem}[Fixed-parameter subcubic perfect-matching search]
\label{thm:subcubic}
For the integers \(a\ge1000\) and fixed reliability parameters
\(0\le\sigma\le a+1\), the construction defines a uniform
randomized algorithm on loopless subcubic graphs with
identified edge copies and \(N\le S\) vertices.
Every computation halts with a checked perfect matching,
a negative answer, or a suspension.
It has no false positive.
For each fixed input with a perfect matching, the
probability of a negative answer is at most
\begin{equation}
\label{phase:pm-error}
                            2^{-\sigma}S^{-5}.
\end{equation}
This statement bounds erroneous answers, not suspensions,
and holds conditionally for adaptive invocations with
fresh randomness and valid prerequisites.

For every sufficiently large fixed \(a\) and fixed
admissible \(\sigma\), there is no suspension on any path
at all sufficiently large \(S\).  The total time on every
path is then
\begin{equation}
\label{phase:pm-time}
                              S^{1+O(1/a)+o_a(1)},
\end{equation}
where the exponent coefficients in \(O(1/a)\) are
absolute.  All initialization, physical data movement,
arithmetic, random bits, and explicit witness writing
are included.
\end{theorem}

\begin{proof}
Settle trivial cases directly; an odd number of vertices
has no perfect matching.  Start with the empty matching
and run the phase construction.  When
\(K\ge200C\), a completed phase gains at least
\(K/(200C)\).  Below that threshold, at most \(200C\)
further phases are needed, each gaining at least one.
Thus there are \(O_a(\log S)\) completed phases, and
at most one final negative or suspended phase.
Each has the bound established above.
Initialization, phase-boundary materialization, and the
final edge-identifier and vertex-coverage check cost
only the already allowed near-linear work.
Lemma~\ref{phase:guard-adequacy} makes
\eqref{phase:pm-time} an every-path bound and excludes
suspension in the stated asymptotic range.
On arbitrary checked sizes, the explicit iteration,
operation, representation, and sampling caps still give
halting, with suspension permitted.

On a PM instance, in the absence of a stochastic failed
guarantee or suspension, valid pair actions and
projections preserve the invariants.
Lemma~\ref{phase:stall} prevents a \(\mathrm{LOW}\)
negative answer before the phase target.  The objective
step bound prevents an infinite sequence of priced
steps, so each phase reaches its gains and the algorithm
eventually constructs a perfect matching.
Checking the proposed edge copies and coverage excludes
false positives on every input, including no-instances.

For the probability bound, count invocations using the
loose guards rather than conditioning on later
noninterruption.  There are fewer than \(S\) nontrivial
phases, at most \(S^{20}\) main actions per phase, and
only logarithmically many threshold tests and retrieval
tests per projection.  On these prefixes,
\[
                      \beta\le CN+L|\mathcal A|,
\]
so its search range has logarithmic bit length.
The total number of pair queries and priced-step
instances is, conservatively, less than \(S^{50}\).
Each reached pair query has the conditional error
allowance in Lemma~\ref{match:pair-action}; each reached
engine instance has the conditional allowance of
Lemma~\ref{lem:priced-steps}.  Both are at most
\(2^{-\sigma}S^{-60}\).
Union over the reached invocations gives at most
\(2^{-\sigma}S^{-10}\), which is smaller than
\eqref{phase:pm-error}.

The conditional bounds are applied when the invocation's
input and prerequisites are fixed by its past history,
and its new bits are fresh.  They are not conditioned on
whether the branch later suspends or wins a scheduling
race.  An additional deterministic suspension produces
no negative answer by itself.  This proves
\eqref{phase:pm-error} also at checked sizes where
suspension remains possible, and completes the theorem.
\end{proof}
